% chapter_07_GenJet.tex -- chapter 07 of the JitJet book: the generator jets (process GENJET).
% Every number comes from test/07_GenJet/out/ (genJetAnalyzer_07_GenJet.log, ledger_07_GenJet.txt, the PNG pages) and
% from the FWLite extras listed in doc/chapter/chapter_07_GenJet.md, section 4.
\providecommand{\antikt}{anti-$k_{\mathrm{T}}$}
\providecommand{\dR}{\ensuremath{\Delta R}}

\chapter{07 -- GenJet: the 634 stable particles are clustered into \antikt{} jets}
\label{ch:07_GenJet}

\section{What this chapter is about}

Up to now the event has been a list of particles. A detector does not measure the $b$ quark of
chapter~\ref{ch:01_ME}, nor the $B^-$ meson of chapter~\ref{ch:05_Had}: it measures a spray of hadrons,
photons and leptons, and the only object that can be compared between theory and experiment is a
\emph{jet}, the sum of the particles inside a cone-like region defined by an algorithm. This chapter runs
that algorithm, \antikt{} as implemented in FastJet 3.4.1, on the 634 stable particles of the full generator
event and produces four collections of \emph{generator jets} (\code{reco::GenJet}): radius $R=0.4$ and
$R=0.8$, each once with all stable particles and once without the neutrinos. The clustering is done as a
separate \code{cmsRun} step, process \code{GENJET}, scheduled by \code{cmsDriver.py} as a \code{USER} step
with the sequence \code{genJetSequence} of \path{JitJetSW/Journey/python/genJet_cff.py}, with the jet
\pt{} threshold lowered to $3\GeV$ so that even the soft jets of the underlying event are visible.
The chapter explains the algorithm, the input selection, and reads the result for the one event: the $b$ jet
of $117.8\GeV$, the $\bar b$ jet of $115.9\GeV$, their balance, and why the $b$ jet is 28\,\% above
the quark it came from but 13\,\% below it once the neutrinos are left out.

\begin{plainwords}[A jet]
A jet is not a particle but a recipe: take the list of final-state particles, apply a clustering rule with a
radius parameter $R$, and call each resulting group a jet, with four-momentum equal to the vector sum of its
constituents. The same recipe applied to the stable particles of the generator gives a \emph{generator jet}
(this chapter), applied to the reconstructed particles of the detector gives a \emph{reconstructed jet}, and
the ratio of the two, the \emph{response}, is what the jet energy corrections (JEC) measure and undo. The
generator jet is therefore the reference of everything that follows in the detector chapters.
\end{plainwords}

\section{From the previous chapter}

The input is \path{test/06_Decay/out/step1_GEN_06_Decay.root}, the complete generator event of
chapter~\ref{ch:06_Decay}: 1846 entries in \code{genParticles}, of which 634 are stable (status~1) and carry
the $13\,600\GeV$ of the two protons, $528.9\GeV$ of scalar \pt{}.\footnote{All numbers of this chapter are
in \path{test/07_GenJet/out/genJetAnalyzer_07_GenJet.log} and \path{ledger_07_GenJet.txt}; the few extra
ones (the origin of the jet constituents, the unclustered particles by region) were read from the output ROOT
file with FWLite and are listed in the companion.} Nothing in that list changes here: the \code{GENJET} step
reads the file, adds the jet collections, and copies \code{genParticles} unchanged under its \code{GEN}
process name. The GEN step of chapter~\ref{ch:06_Decay} had in fact already clustered jets (the standard
\code{pgen} sequence contains \code{genJetMET}); those are dropped with
\code{--outputCommands "drop *\_ak4GenJets*\_*\_GEN,drop *\_ak8GenJets*\_*\_GEN"}, so that the output holds
exactly one set of jets, the one made and configured in this chapter.

The two objects to follow are the $b$ and $\bar b$ quarks. The hard-process $b$ (status~23,
$\pt=113.2\GeV$, $\eta=-1.44$) is pushed to $122.8\GeV$ by the recoil of the initial-state radiation
(status~44), radiates twice in the final-state shower (status~52: $106.0$, then $91.6\GeV$), receives its
primordial $k_{\mathrm{T}}$ (status~62, $91.7\GeV$) and is copied as the end of a colour string (status~71,
$\pt=91.74\GeV$, $\eta=-1.57$, $\phi=-2.08$). This status-71 copy is the quark the analyzer compares the
jets with. It fragments into the $B^-$ of $94.3\GeV$, which decays semileptonically,
$B^-\to D_0^{*0}\,e^-\,\bar\nu_e$, with the antineutrino carrying $37.5\GeV$ of \pt{} ($95.5\GeV$ of energy)
and nine stable decay products in all. The $\bar b$ goes from $113.2$ to $116.4\GeV$ (status~71,
$\eta=-3.75$, $\phi=1.01$), fragments into a $B^{*+}$ of $114.2\GeV$ and ends as 16 stable particles, none
of them a neutrino. At parton level the pair is already unbalanced: $|\vec p_{\mathrm T}(b)+\vec p_{\mathrm
T}(\bar b)|=25.1\GeV$ for the status-62 quarks and $20.3\GeV$ for the two $B$ hadrons, against $0$ for the
hard process. Jets exist to recover what the radiation took away.

\section{The numbers of the event}

\begin{plainwords}[The \antikt{} algorithm]
FastJet computes for every pair of particles $d_{ij}=\min(1/p_{\mathrm{T}i}^2,\,1/p_{\mathrm{T}j}^2)\,
\Delta R_{ij}^2/R^2$ with $\Delta R_{ij}^2=(\eta_i-\eta_j)^2+(\phi_i-\phi_j)^2$, and for every particle
$d_{iB}=1/p_{\mathrm{T}i}^2$ (the distance to the beam). The smallest of all these numbers decides: if it is a
$d_{ij}$, $i$ and $j$ are merged (four-momenta added) and the list is recomputed; if it is a $d_{iB}$, $i$ is
declared a jet and removed. The inverse powers of \pt{} make the hardest particle the first to act: it
swallows everything within \dR$<R$ around it before any soft particle can merge with another soft one, so a
hard jet is a circle of radius $R$ around its axis, as in the display, and the soft background does not move
its axis. Two properties make the definition usable in perturbative QCD: a soft particle added anywhere does
not change the hard jets (\emph{infrared safety}), and a particle split into two collinear ones gives the
same jets (\emph{collinear safety}); a cone algorithm seeded on the hardest particle has neither. The jet
\pt{} is the \pt{} of the vector sum, slightly smaller than the scalar sum of its constituents.
\end{plainwords}

\textbf{The producer and its input.} Each collection is one \code{FastjetJetProducer} module
(\path{RecoJets/JetProducers/python/ak4GenJets_cfi.py}: \code{jetAlgorithm} \code{"AntiKt"}, \code{rParam} 0.4;
\path{ak8GenJets_cfi.py} clones it with \code{rParam} 0.8) that reads the parameter set
\code{GenJetParameters} (\path{GenJetParameters_cfi.py}): \code{jetType} \code{'GenJet'},
\code{jetPtMin} \code{3.0}, \code{inputEtMin} \code{0.0}, \code{inputEMin} \code{0.0}, no area, no $\rho$ subtraction, a
deterministic FastJet seed. \code{jetPtMin} is the only threshold: jets below it are not written. The script
sets it explicitly to the same $3\GeV$ for the four producers, through \code{--customise\_commands}, so that it
can be changed with \code{JETPTMIN}. The input is not \code{genParticles} itself but a selection made by
\code{InputGenJetsParticleSelector} (\path{RecoJets/Configuration/python/GenJetParticles_cff.py}): it keeps
every status-1 particle except those whose $|\mathrm{pdgId}|$ is in \code{ignoreParticleIDs}. For
\code{genParticlesForJets} that list contains only hypothetical invisible particles: 1000022 (the lightest
neutralino), 1000012/14/16 and 2000012/14/16 (left- and right-handed sneutrinos), 1000039 (gravitino),
5100039 (Kaluza--Klein graviton), 4000012/14/16 (excited neutrinos), 9900012/14/16 (heavy Majorana
neutrinos) and 39 (graviton). None of them occurs in a Pythia8 \bbbar{} event, so all 634 stable particles
are offered to the clustering. \code{genParticlesForJetsNoNu} adds 12, 14, 16, the three neutrinos: 633
particles. The other switches are off: \code{excludeResonances = False} (no removal of $W/Z/H$ decay
products), \code{tausAsJets = False}, \code{partonicFinalState = False}.

\begin{table}[htb]
\centering\small
\caption{The $b$ and $\bar b$ jets in the three collections drawn (\code{genJetAnalyzer\_07\_GenJet.log}).
The quark is the status-71 copy ($\pt=91.74$ and $116.36\GeV$); \dR{} is the distance between jet axis
and quark; the net \pt{} and $\Delta\phi$ are those of the pair; $m_{jj}$ is read from the page header
(\code{genJetAnalyzer\_07\_GenJet\_page-6/7/8.png}).}
\label{tab:07_bjets}
\begin{tabular}{llrrrrrrr}
\toprule
collection & jet & \pt{} [GeV] & $\eta$ & $\phi$ & $n$ & \dR{} & $\pt^{\mathrm{jet}}/\pt^{q}$ & mass [GeV] \\
\midrule
\code{ak4GenJets}     & $b$       & 117.79 & $-1.59$ & $-2.08$ & 24 & 0.019 & 1.284 & 16.7 \\
                      & $\bar b$  & 115.94 & $-3.75$ & $1.01$  & 18 & 0.004 & 0.996 & 7.3 \\
                      & pair      & \multicolumn{7}{l}{net $\pt=6.06\GeV$, $\Delta\phi=3.09$, $\Sigma\pt=233.7\GeV$, $m_{jj}=386.8\GeV$} \\
\addlinespace
\code{ak4GenJetsNoNu} & $b$       & 79.82  & $-1.59$ & $-2.08$ & 22 & 0.017 & 0.870 & 13.4 \\
                      & $\bar b$  & 115.94 & $-3.75$ & $1.01$  & 18 & 0.004 & 0.996 & 7.3 \\
                      & pair      & \multicolumn{7}{l}{net $\pt=36.53\GeV$, $\Delta\phi=3.09$, $\Sigma\pt=195.8\GeV$, $m_{jj}=319.7\GeV$} \\
\addlinespace
\code{ak8GenJets}     & $b$       & 141.10 & $-1.65$ & $-2.05$ & 55 & 0.088 & 1.538 & 47.3 \\
                      & $\bar b$  & 121.01 & $-3.75$ & $0.99$  & 32 & 0.022 & 1.040 & 19.0 \\
                      & pair      & \multicolumn{7}{l}{net $\pt=24.03\GeV$, $\Delta\phi=3.04$, $\Sigma\pt=262.1\GeV$, $m_{jj}=439.5\GeV$} \\
\bottomrule
\end{tabular}
\end{table}

\textbf{The $b$ jet and its response.} With $R=0.4$ the $b$ quark is matched to jet~[0], $\pt=117.79\GeV$,
24 constituents, axis $0.019$ away from the quark direction. Its response with respect to the status-71 quark
is $117.79/91.74=1.284$. A response above one is not an excess of energy, it is the bookkeeping of the
shower: the quark used as reference has already radiated. Measured against the hard-process $b$ of
$113.2\GeV$ the same jet has a response of $1.04$. The constituents tell where the \pt{} comes from. Nine of
the 24 are the decay products of the $B^-$ and carry $94.3\GeV$ (the neutrino alone $37.5\GeV$). Seven more,
$17.5\GeV$, are hadrons from the same colour string as the $b$ quark, the string formed by the $b$, three
gluons and an $\bar s$ quark of the shower: this is the radiated energy coming back. The remaining eight,
$6.4\GeV$, come from other strings, those of the multiparton interactions and of the beam remnants: the
underlying event that happens to fall inside the circle. The $\bar b$ jet at $\eta=-3.75$ is cleaner: 16 of
its 18 constituents are the $B^{*+}$ decay products ($114.2\GeV$), one $K^-$ of $1.6\GeV$ comes from the
$\bar b$ string and one photon of $0.15\GeV$ from elsewhere; response $0.996$.

\begin{plainwords}[Why 1.28 with neutrinos and 0.87 without]
The $B^-$ decayed semileptonically and its antineutrino carries $37.5\GeV$ of the jet's $117.8\GeV$.
\code{ak4GenJetsNoNu} drops the three neutrino species before clustering: the $b$ jet becomes
$79.82\GeV$ (response $0.870$), with 22 constituents instead of 24 because, with the neutrino gone, the axis
moved and a $0.5\GeV$ pion at the edge now belongs to the neighbouring jet~[4] ($11.89\to12.40\GeV$, 11 to 12
constituents). The $\bar b$ jet, without a neutrino, is identical in both collections. The dijet balance
follows: with neutrinos the pair has a net \pt{} of $6.1\GeV$ and $\Delta\phi=3.09$, the best balance of the
whole journey ($25.1\GeV$ at quark level); without, $36.5\GeV$ is missing in the direction of the $b$ jet,
exactly the signature a detector would call missing transverse energy. Which of the two is the right reference
depends on the question: the NoNu jets are what a calorimeter can see at best and are the reference of the
jet energy corrections; the full jets are the parton-level truth.
\end{plainwords}

\textbf{$R=0.4$ against $R=0.8$.} Doubling the radius quadruples the area. The $b$ jet of \code{ak8GenJets}
has $141.10\GeV$ and 55 constituents, response $1.538$ (1.25 against the hard-process quark), mass
$47.3\GeV$ instead of $16.7$; its axis is $0.088$ from the quark. The additions are $19.8\GeV$ of further
hadrons of the $b$ string (8 particles, in all 17 of the 55 constituents and $114.2\GeV$ belong to the $b$
string) and $28.5\GeV$ from 38 particles of other strings. The $\bar b$ jet gains $5.1\GeV$ ($121.01\GeV$,
32 constituents). The pair is now less balanced, net $\pt=24.0\GeV$, $\Delta\phi=3.04$: the large cone
collects the underlying event unevenly. In \code{ak8GenJets} 488 of the 634 stable particles are in a jet and
28 jets are found; jets [2] to [7] have $10$ to $29\GeV$ and 17 to 37 constituents each, merged soft activity
that \code{ak4GenJets} splits into several $5$ to $14\GeV$ jets. The small radius is the one used for most
of CMS physics, with the underlying event and (later) the pileup as the price of a large one.

\textbf{The other jets and the unclustered particles.} \code{ak4GenJets} finds 27 jets above $3\GeV$. Besides
the two $b$ jets, jets~[2] to [5] ($13.8$, $13.3$, $11.9$, $11.4\GeV$, 10 to 12 constituents) are the
hadronised radiation and multiparton interactions of chapters~\ref{ch:02_PS} and \ref{ch:04_MPI}; the
remaining 21 have $3.2$ to $7.1\GeV$. In all, 239 of the 634 stable particles are inside a jet and the jets
sum to $382.8\GeV$ of \pt{}. The 395 unclustered particles carry $144.6\GeV$ of scalar \pt{}: 298 of them are
below $0.5\GeV$, 78 between $0.5$ and $1\GeV$, only two above $2\GeV$ (the hardest $2.3\GeV$), and 38 sit
beyond $|\eta|=6$, on the border of the display. They carry, however, $9615\GeV$ of energy, almost all of it
beam remnants along the beam pipe, which is why the jet energy ($3985\GeV$) is a small part of the event
energy. By region, the 27 \code{ak4GenJets} are 17 at $\abseta<2.5$ ($217.9\GeV$), 4 at $2.5<\abseta<3.0$
($18.5\GeV$) and 6 at $3.0<\abseta<5.2$ ($146.3\GeV$, among them the $\bar b$ jet); none is beyond $5.2$,
while \code{ak8GenJets} has three jets there ($16.1\GeV$, up to $\eta=6.1$). Generator jets have no $\eta$
limit; the detector will set one.

\section{Ledger and display}

\begin{table}[H]
\centering\small
\caption{Ledger of the GenJet stage (\code{ledger\_07\_GenJet.txt}; the unclustered and per-region lines from
the FWLite extras of the companion). The stage adds nothing to and removes nothing from the event: it
partitions the 634 stable particles into jets and a remainder. The jet \pt{} is the vector sum per jet, the
particle sums are scalar.}
\label{tab:07_ledger}
\begin{tabular}{lrrrrrr}
\toprule
& \multicolumn{3}{c}{before: stable particles} & \multicolumn{3}{c}{after: \code{ak4GenJets}} \\
\cmidrule(lr){2-4}\cmidrule(lr){5-7}
 & $n$ & $E$ [GeV] & \pt{} [GeV] & jets & $n_{\mathrm{const}}$ & $\Sigma\pt$ [GeV] \\
\midrule
all, in a jet           & 634 & 13600.0 & 528.9 & 27 & 239 & 382.8 \\
all, not in a jet       &     &         &       & -- & 395 & 144.6 \\
$\abseta<2.5$           & 316 & 733.8   & 278.7 & 17 &     & 217.9 \\
$2.5<\abseta<3.0$       & 55  & 193.1   & 24.6  & 4  &     & 18.5 \\
$3.0<\abseta<5.2$       & 189 & 5019.6  & 197.4 & 6  &     & 146.3 \\
$\abseta>5.2$           & 74  & 7653.5  & 28.1  & 0  &     & 0 \\
neutrinos               & 1   & 95.5    & 37.5  & \multicolumn{3}{l}{in \code{ak4GenJets}, not in \code{NoNu}} \\
\midrule
\code{ak4GenJets}       &     &         &       & 27 & 239 & 382.8 \ ($\Sigma E=3985.2$) \\
\code{ak4GenJetsNoNu}   &     &         &       & 27 & 238 & 345.3 \ ($\Sigma E=3889.7$) \\
\code{ak8GenJets}       &     &         &       & 28 & 488 & 473.1 \ ($\Sigma E=7613.4$) \\
\bottomrule
\end{tabular}

\medskip
\begin{tabular}{lrrrrrr}
\toprule
$b+\bar b$ pair [GeV] & hard process & status 62 & $B$ hadrons & \code{ak4GenJets} & \code{NoNu} & \code{ak8GenJets} \\
\midrule
scalar sum \pt{} & 226.4 & 208.1 & 208.5 & 233.7 & 195.8 & 262.1 \\
net \pt{}        & 0.0   & 25.1  & 20.3  & 6.1   & 36.5  & 24.0 \\
\bottomrule
\end{tabular}
\end{table}

\begin{figure}[H]
\centering
\includegraphics[width=0.8\textwidth]{07_GenJet/out/genJetAnalyzer_07_GenJet_page-6.png}
\caption{\code{ak4GenJets}: the 27 \antikt{} $R=0.4$ jets with $\pt>3\GeV$ of the event, over the 634 stable
particles (marker size $\sim\log\pt$). The $b$ jet [0] ($117.8\GeV$, red) and the $\bar b$ jet [1]
($115.9\GeV$, blue) are matched to the status-71 quarks (stars). Pages 1--5 of the same PDF are the five
generator stages of chapters~\ref{ch:01_ME}--\ref{ch:06_Decay}.
(\code{test/07\_GenJet/out/genJetAnalyzer\_07\_GenJet\_page-6.png})}
\label{fig:07_ak4}
\end{figure}

Figure~\ref{fig:07_ak4} is the $\eta$--$\phi$ display of the event at this stage: every jet is a circle of
radius $R$ around its axis, its constituents in the jet colour, the unclustered stable particles in grey, the
$b$ and $\bar b$ jets in red and blue with the quark direction as an open star. The two $b$ jets are
back-to-back in $\phi$ ($\Delta\phi=3.09$) and far apart in $\eta$ ($-1.59$ and $-3.75$), the picture of a
$\qqbar\to\bbbar$ event with a very asymmetric $x_1=0.39$, $x_2=0.0022$ of chapter~\ref{ch:01_ME}. The
circles of the 25 other jets, mostly in the central region, are the hadronised shower and MPI. The grey
particles between the circles and the 38 crosses on the $|\eta|=6$ border are the beam remnants.
Figure~\ref{fig:07_nonu_ak8} shows the same event without neutrinos (the $b$ jet shrinks to $79.8\GeV$,
nothing else moves) and with $R=0.8$ (the circles overlap, 488 particles are clustered, the display fills
up).


\begin{figure}[htb]
\centering
\begin{subfigure}[t]{0.495\textwidth}
\includegraphics[width=\textwidth]{07_GenJet/out/genJetAnalyzer_07_GenJet_page-7.png}
\caption{\code{ak4GenJetsNoNu}: without the neutrino the $b$ jet is $79.8\GeV$.}
\end{subfigure}\hfill
\begin{subfigure}[t]{0.495\textwidth}
\includegraphics[width=\textwidth]{07_GenJet/out/genJetAnalyzer_07_GenJet_page-8.png}
\caption{\code{ak8GenJets}: $R=0.8$, the $b$ jet is $141.1\GeV$ with 55 constituents.}
\end{subfigure}
\caption{The same event in the two other collections
(\code{genJetAnalyzer\_07\_GenJet\_page-7.png}, \code{page-8.png}).}
\label{fig:07_nonu_ak8}
\end{figure}

\FloatBarrier
\section{Open problems for Phase-2}

The generator jets are the reference of the jet energy scale, so what is frozen here is frozen for the whole
HL-LHC programme. A look into \code{CMSSW\_15\_0\_5} gives the following.

\textbf{The definitions have not moved since Run 2.} \path{RecoJets/JetProducers/python/GenJetParameters_cfi.py}
sets \code{jetPtMin} \code{3.0}, \code{inputEtMin} \code{0.0}, \code{inputEMin} \code{0.0}, \code{doAreaFastjet}
\code{False}, \code{doRhoFastjet} \code{False}; \path{ak4GenJets_cfi.py} and \path{ak8GenJets_cfi.py} fix \code{rParam} to 0.4
and 0.8. No file in \path{RecoJets/Configuration/python} or \path{RecoJets/JetProducers/python} carries a
\code{phase2} era modifier for the generator jets (the only \code{phase2} strings in \code{RecoJets} are the
heavy-ion event content, \path{RecoJets_EventContent_cff.py}, and \path{CaloTowersES_cfi.py}); the Phase-2
GEN step of \path{doc/tmp/Gen4Phase2.md} runs byte-identical producers. The thresholds that matter
downstream are also unchanged: miniAOD keeps \code{slimmedGenJets} from \code{ak4GenJetsNoNu} without a
cut and \code{slimmedGenJetsAK8} from \code{ak8GenJetsNoNu} with \code{pt > 80}
(\path{PhysicsTools/PatAlgos/python/slimming/slimmedGenJets_cfi.py}); nanoAOD writes \code{GenJet} with
\code{pt > 10} and \code{GenJetAK8} with \code{pt > 100} (\path{PhysicsTools/NanoAOD/python/jetMC_cff.py}).
Note that the compiled default of \code{jetPtMin} in \code{VirtualJetProducer::fillDescriptions} is $5\GeV$
while the \code{cfi} says $3\GeV$: the \code{cfi} wins, but a module built without it would cut differently.
None of this is wrong for Phase-2, but none of it was revisited for it either.

\textbf{The exclusion lists are the 2007 lists.} The file
\path{RecoJets/Configuration/python/GenJetParticles_cff.py} is headed ``ShR 27 Mar 07'' and its \code{ignoreParticleIDs} are the invisible BSM particles listed above; the
only later additions are the \code{NoNu} (12, 14, 16) and \code{NoMuNoNu} (12, 13, 14, 16) clones. There is
no pileup-robust generator-jet definition and there need not be one: \code{genParticles} holds only the
signal collision, so a generator jet is pileup-free by construction (chapter~\ref{ch:08_PU} shows what happens
when the pileup particles are added to the clustering on purpose). The pileup robustness is entirely on the
reconstruction side, and that side has been retuned for Phase-2: \path{CommonTools/PileupAlgos/python/Puppi_cff.py}
modifies \code{puppi} under \code{phase2\_common} with \code{EtaMinUseDeltaZ = 4.0}, \code{DeltaZCut = 0.1},
\code{MinNeutralPt = 0.2} up to $\abseta=3.5$ and $2.0$ beyond, with the $\eta$ regions $0$--$2.5$,
$2.5$--$3.5$, $3.5$--$10$. The reference against which that tuning is judged remains the unchanged
\code{ak4GenJetsNoNu}. One open point for a $b$-rich analysis is the one visible in this event: the
\code{NoNu} reference removes the $37.5\GeV$ neutrino and makes the $b$ jet response $0.87$, so the flavour
dependence of the corrections, measured against this reference, is as large at the HL-LHC as it is now.

\textbf{The $\eta$ reach.} Generator jets have no $\eta$ cut; in this event \code{ak4GenJets} reach
$\eta=-4.70$ and \code{ak8GenJets} $\eta=6.12$. The Run-3 detector of this book sees jets up to
$\abseta\approx5.2$ with the HF, with the region $2.5$--$3.0$ covered by calorimeters only; the Phase-2 layout
in the release puts the tracker to $\abseta=4.0$ (the \code{EtaMinUseDeltaZ = 4.0} of Puppi above), the
HGCal at $1.5<\abseta<3.0$ and the HF at $3.0<\abseta<5.5$, as spelled out by the region boundaries of the
Phase-2 Level-1 particle flow, \path{L1Trigger/Phase2L1ParticleFlow/python/l1ctLayer1_cff.py}
(\code{etaBoundaries = (-1.5, 1.5)}, \code{(1.5, 3.0)}, \code{(3.0, 5.5)}). In the ledger of this event the
$2.5<\abseta<3.0$ region holds 4 \code{ak4GenJets} ($18.5\GeV$) and $3.0<\abseta<5.2$ holds 6
($146.3\GeV$, the $\bar b$ jet among them): with the HGCal and the extended tracker the $\bar b$ jet of
this event, at $\eta=-3.75$, is still an HF jet, and its $B^{*+}$ decay products at $\abseta>3$ are still
untracked. Whether the generator-jet validation follows the new boundaries could not be established:
\path{Validation/RecoJets/python} contains no \code{phase2} modifier.

\textbf{The jet energy corrections of Phase-2.} The chain that turns a reconstructed jet into a
generator-jet-like one is a sequence of \code{JetCorrectorParametersCollection} payloads in the global tag.
For the Phase-2 tag of this release, \path{auto:phase2_realistic} $=$ \path{150X_mcRun4_realistic_v1}
(\path{Configuration/AlCa/python/autoCond.py}), \code{conddb list} gives: \code{AK4PF}, \code{AK8PF} and
\code{AK8PFchs} from \path{PhaseIIFall17_V5_MC}, a 2017 Phase-2 study; \code{AK4PFPuppi},
\code{AK4PFchs} and \code{AK8PFPuppi}, the collections actually used offline, from
\path{Winter22Run3_V1_MC}, i.e.\ a Run-3 payload applied to the Phase-2 detector; only the HLT one,
\code{AK4PFPuppiHLT}, is recent and Phase-2 (\path{Phase2Spring23ticlv4_MC_v1}). The jet resolution
payloads of the same tag are all \path{JR_Spring16_25nsV6_MC}, from 2016. For comparison the 2024 tag of
this book, \path{150X_mcRun3_2024_realistic_v2}, has \path{Winter24_AK4PFPuppi_offline_v1} and
\path{JR_Summer19UL18_JRV2}. All of them take \code{ak4GenJetsNoNu} (or \code{ak8GenJetsNoNu}) matched
to the reconstructed jet as the truth: the generator jets of this chapter, with their 2007 input list and
$3\GeV$ threshold, are the reference of the HL-LHC jet energy scale until a Phase-2 payload is derived, and
they will have to serve at PU 200 and $14\TeV$ without having been changed for either. What could not be
established from the release: the $\eta$ binning of the Phase-2 payloads (the files live only in the
database, \path{CondFormats/JetMETObjects/data} holds no Run-3 or Phase-2 text file) and whether the matching
radius and \pt{} thresholds used to derive them will be adjusted for the HGCal granularity.
